
\documentclass{article}
\usepackage{colortbl}
\usepackage{makecell}
\usepackage{multirow}
\usepackage{supertabular}

\begin{document}

\newcounter{utterance}

\centering \large Interaction Transcript for game `cryptolect', experiment `standard', episode 7 with BSU{-}SLIM\_\_Qwen3.5{-}9B{-}Base\_\_playpen{-}prm{-}9b{-}best\_of\_n.
\vspace{24pt}

{ \footnotesize  \setcounter{utterance}{1}
\setlength{\tabcolsep}{0pt}
\begin{supertabular}{c@{$\;$}|p{.15\linewidth}@{}p{.15\linewidth}p{.15\linewidth}p{.15\linewidth}p{.15\linewidth}p{.15\linewidth}}
    \# & \multicolumn{2}{c}{Player} && \multicolumn{2}{c}{Game Master} \\
    \hline

    \theutterance \stepcounter{utterance}  
    & & & \multicolumn{4}{p{0.6\linewidth}}{
        \cellcolor[rgb]{0.9,0.9,0.9}{
            \makecell[{{p{\linewidth}}}]{
                \texttt{\tiny{[P1$\langle$GM]}}
                \texttt{You are a field linguist. You have made contact with the only speaker of an unknown} \\
\texttt{language (the "informant"). Nobody has ever written down this language before. Your task} \\
\texttt{is to work out how it works and then translate sentences into it correctly.} \\
\\ 
\texttt{You will study the language by asking the informant to translate English sentences for} \\
\texttt{you. You may only talk about a small fixed world, using exactly these English words:} \\
\\ 
\texttt{Nouns: king, bird, hunter, farmer, dog, hill, spear, lake, road, boat, tree, stone, river, ship} \\
\texttt{Adjectives: big, red, cold, blue, small, tall, old, dark} \\
\texttt{Verbs (write them in this base form): see, chase, feed, find, watch, hold, carry, hear} \\
\texttt{Articles: the (definite), a (indefinite)} \\
\texttt{Number words: two, three, four, five, six (no number word = one / singular)} \\
\texttt{These are ALL the words in play {-} the language has no other words.} \\
\\ 
\texttt{HOW TO WRITE AN ENGLISH SENTENCE} \\
\texttt{Every sentence is one clause with a subject, a verb, and an object, written in this fixed English notation:} \\
\texttt{    SUBJECT  verb  OBJECT} \\
\texttt{Each of SUBJECT and OBJECT is a noun phrase of the form:} \\
\texttt{    {[}the | a{]}  {[}number word{]}  {[}adjective ...{]}  noun} \\
\texttt{{-} 'the' = definite, 'a' = indefinite (singular only); leave the article off for an indefinite plural.} \\
\texttt{{-} a number word (two, three, four, five, six) sets the count; with no number word the noun is singular.} \\
\texttt{{-} zero or more adjectives may come before the noun.} \\
\texttt{Examples of well{-}formed English sentences you may send:} \\
\texttt{    a king see a bird} \\
\texttt{    six kings see a bird} \\
\texttt{    a big king see a bird} \\
\\ 
\texttt{WHAT YOU ALREADY KNOW} \\
\texttt{The informant has already translated these sentences for you. Each line is} \\
\texttt{"<English>  =>  <the language>":} \\
\texttt{    the five lakes carry the big farmer  =>  gugepa fugo fogupo getupopa ti} \\
\texttt{    five hunters watch the two red ships  =>  gisapopa fago te kutopi sukupo ti} \\
\texttt{    a blue road feed the five farmers  =>  gugepopa ti totu tigo lofa} \\
\texttt{    a cold road find three ships  =>  gisapo pi kutofo tigo gugefe} \\
\texttt{    the four hunters find the three big rivers  =>  gesufapopa fugo pi kutofo sukupopa se} \\
\texttt{    the road find three red roads  =>  tigopo fago pi kutofo tigopa} \\
\texttt{    the two boats chase the four kings  =>  logipopa se tuli gasitipopa te} \\
\texttt{    three stones find the two old trees  =>  palepopa giso te kutofo gulupo pi} \\
\texttt{    three farmers watch the six rivers  =>  gesufapopa pufa kutopi gugepo pi} \\
\texttt{    the two rivers watch a tall hill  =>  kugi gole kutopi gesufapopa te} \\
\texttt{    the two hills feed three dark dogs  =>  fapopo lulo pi totu kugipopa te} \\
\texttt{    the five birds carry six rivers  =>  gesufapo pufa fogupo kafepopa ti} \\
\texttt{    the four dark boats hold four dark spears  =>  kasefupo lulo se gile gasitipopa lulo se} \\
\texttt{    the two lakes hear a dark hill  =>  kugi lulo fosu getupopa te} \\
\texttt{    two small birds hear a dog  =>  fapo fosu kafepo lesa te} \\
\texttt{    a hill see a tall river  =>  gesufa gole kake kugi} \\
\\ 
\texttt{HOW THE SESSION WORKS} \\
\texttt{1. INVESTIGATION: You have 5 rounds of questioning. In each round you may ask} \\
\texttt{   the informant to translate up to 4 English sentences (so at most} \\
\texttt{   20 probes in total). The informant answers each one truthfully and} \\
\texttt{   exactly. Use these answers to figure out the words and the grammar (word order, how} \\
\texttt{   number/case/definiteness are marked, how words must agree, and any sound changes).} \\
\texttt{   Choose your probes deliberately {-} test specific hypotheses.} \\
\texttt{2. EXAM: After the investigation rounds, the informant will give you 10 NEW} \\
\texttt{   English sentences. You must translate each one INTO the language. Every answer is} \\
\texttt{   graded by EXACT match against how the informant would say it {-} every word, ending and} \\
\texttt{   space must be right (capitalisation and surrounding spaces are ignored). The exam} \\
\texttt{   sentences will use the same vocabulary but in combinations you have not been shown, so} \\
\texttt{   you must induce the general rules, not just memorise examples.} \\
\\ 
\texttt{RESPONSE FORMAT (for every investigation round)} \\
\texttt{List your probe sentences under a 'PROBES:' line, one per line. A THOUGHTS line is} \\
\texttt{optional {-} keep any reasoning brief so you always finish writing your PROBES:} \\
\texttt{THOUGHTS: <optional; brief>} \\
\texttt{PROBES:} \\
\texttt{1. <an English sentence, written as described above>} \\
\texttt{2. <another English sentence>} \\
\texttt{(up to 4 lines; at least one)} \\
\\ 
\texttt{Begin your first round of probes now.} \\
            }
        }
    }
     \\ \\

    \theutterance \stepcounter{utterance}  
    & \multicolumn{4}{p{0.6\linewidth}}{
        \cellcolor[rgb]{0.561,0.737,0.561}{
            \makecell[{{p{\linewidth}}}]{
                \texttt{\tiny{[P1$\rangle$GM]}}
                \texttt{THOUGHTS: I need to start mapping the language by identifying the base words for the known nouns and verbs, and understanding how they change based on definiteness (the/a), number (singular/plural), and adjectives. I also need to determine the word order (Subject{-}Verb{-}Object seems likely from the English input structure) and whether there are suffixes or prefixes indicating case or agreement.} \\
\\ 
\texttt{Looking at the first few examples:} \\
\texttt{1. `the five lakes carry the big farmer` => `gugepa fugo fogupo getupopa ti`} \\
\texttt{   {-} Noun 'lakes' (plural, definite) seems to be related to `fugo`.} \\
\texttt{   {-} Noun 'farmer' (singular, definite) seems to be related to `ti`.} \\
\texttt{   {-} Verb 'carry' seems to be `getupopa`.} \\
\texttt{   {-} Adjective 'big' seems to be `fug{-}` (attached to `fugo`? `fugo` vs `fugupo`?). Wait, `fugo` is the base for 'lake'? Let's check other examples.} \\
\texttt{   {-} `five lakes`: `five` {-}> `fugo`? Or does `fugo` mean 'lakes'?} \\
\texttt{   {-} Let's look at `the four hunters`: `gesufapopa fugo`. Here 'hunters' (plural, definite) is `gesufapopa`, and `fugo` appears again. In example 1, `fugo` was 'lakes'. This suggests `fugo` might be a stem for both 'hunter' and 'lake' OR the verb agreement is causing confusion.} \\
\texttt{   {-} Let's re{-}examine: `the five lakes carry...` (Subject: lakes). `...carry` (Verb). `...the big farmer` (Object: farmer).} \\
\texttt{   {-} Output: `gugepa fugo fogupo getupopa ti`} \\
\texttt{   {-} If SVO: Subject=`gugepa fugo fogupo`, Verb=`getupopa`, Object=`ti`.} \\
\texttt{   {-} Or VSO? `getupopa gugepa fugo fogupo ti`? No, the input is SVO.} \\
\texttt{   {-} Hypothesis: The language might be SVO.} \\
\texttt{   {-} Let's check `five hunters watch the two red ships` => `gisapopa fago te kutopi sukupo ti`.} \\
\texttt{     {-} Subject: `five hunters`. Output segment: `gisapopa fago te`.} \\
\texttt{     {-} Verb: `watch`. Output segment: `sukupo`.} \\
\texttt{     {-} Object: `the two red ships`. Output segment: `ti`.} \\
\texttt{     {-} So: `gisapopa fago te` (Subject) `sukupo` (Verb) `ti` (Object).} \\
\texttt{     {-} Compare to Ex 1: `gugepa fugo fogupo` (Subject) `getupopa` (Verb) `ti` (Object).} \\
\texttt{     {-} In Ex 1, Object 'big farmer' is `ti`. In Ex 2, Object 'two red ships' is `ti`.} \\
\texttt{     {-} This strongly suggests `ti` is the word for 'farmer' AND 'ships'? That seems unlikely unless the lexicon is very small/crossed.} \\
\texttt{     {-} Let's check other examples for `ti`.} \\
\texttt{     {-} `a blue road feed the five farmers` => `gugepopa ti totu tigo lofa`.} \\
\texttt{       {-} Subject: `gugepopa ti`. Verb: `totu`. Object: `tigo lofa`.} \\
\texttt{       {-} Here `ti` is part of the Subject ("blue road"). So 'road' is `ti`?} \\
\texttt{       {-} But in Ex 1 \& 2, `ti` was the object ('farmer', 'ships').} \\
\texttt{       {-} This is confusing. Maybe the word order is different? Or maybe `ti` is a verb?} \\
\texttt{       {-} Let's re{-}evaluate Ex 1: `the five lakes carry the big farmer`.} \\
\texttt{         {-} If `ti` is 'farmer', then Subject is `gugepa fugo fogupo` ('five lakes'?).} \\
\texttt{         {-} If `ti` is 'river'? Example 15: `two small birds hear a dog` => `fapo fosu kafepo lesa te`. 'dog' is `lesa`? 'birds' is `kafepo`?} \\
\texttt{         {-} Example 17: `a hill see a tall river` => `gesufa gole kake kugi`. 'hill' is `gesufa`? 'river' is `kugi`?} \\
\texttt{       {-} Let's try to isolate roots.} \\
\texttt{       {-} Ex 17: `a hill` {-}> `gesufa`. `a tall river` {-}> `kake kugi`. So 'river' is likely `kugi`.} \\
\texttt{       {-} Ex} \\
            }
        }
    }
    & & \\ \\

    \theutterance \stepcounter{utterance}  
    & & & \multicolumn{2}{p{0.3\linewidth}}{
        \cellcolor[rgb]{0.9,0.9,0.9}{
            \makecell[{{p{\linewidth}}}]{
                \texttt{\tiny{[GM$|$GM]}}
                \texttt{could not find any probe sentences; put them under a 'PROBES:' line, one per line} \\
            }
        }
    }
    & & \\ \\

    \theutterance \stepcounter{utterance}  
    & & & \multicolumn{4}{p{0.6\linewidth}}{
        \cellcolor[rgb]{0.9,0.9,0.9}{
            \makecell[{{p{\linewidth}}}]{
                \texttt{\tiny{[P1$\langle$GM]}}
                \texttt{Your last response could not be processed: could not find any probe sentences; put them under a 'PROBES:' line, one per line.} \\
\texttt{List your probes (a brief THOUGHTS line is optional):} \\
\texttt{PROBES:} \\
\texttt{1. <an English sentence>} \\
\texttt{(up to 4 probe lines; at least one)} \\
            }
        }
    }
     \\ \\

    \theutterance \stepcounter{utterance}  
    & \multicolumn{4}{p{0.6\linewidth}}{
        \cellcolor[rgb]{0.561,0.737,0.561}{
            \makecell[{{p{\linewidth}}}]{
                \texttt{\tiny{[P1$\rangle$GM]}}
                \texttt{THOUGHTS: I need to identify the base words for the nouns and verbs first. From the data provided:} \\
\texttt{{-} "the five lakes carry..." {-}> Subject has "lakes". Output starts with "gugepa fugo fogupo".} \\
\texttt{{-} "five hunters watch..." {-}> Subject has "hunters". Output starts with "gisapopa fago te".} \\
\texttt{{-} "a blue road feed..." {-}> Subject has "road". Output starts with "gugepopa ti".} \\
\texttt{{-} "the road find..." {-}> Subject has "road". Output starts with "tigopo fago pi".} \\
\texttt{Comparing "road" examples: "gugepopa ti" vs "tigopo fago pi". The adjective "blue" is "tigo" in the first, but "fago" in the second? Wait, "red" was "fago" in "the two red ships". Let's re{-}check "red".} \\
\texttt{{-} "the two red ships" {-}> "fago te kutopi". Here "fago" modifies "kutopi". So "red" is "fago".} \\
\texttt{{-} In "the road find three red roads", "red" modifies "roads". Output is "tigopo fago pi kutofo". So "fago" is the adjective.} \\
\texttt{{-} But in "a blue road feed...", "blue" is "tigo". So "blue" = "tigo".} \\
\texttt{{-} Now looking at "road": In "a blue road", it is "ti". In "the road", it is "tigopo".} \\
\texttt{{-} Hypothesis: "ti" is the base word for "road". "tigopo" might be "the" + "road" or "road" + plural marker? Or maybe "road" is "ti" and the suffix changes for definiteness/number.} \\
\texttt{{-} Let's look at verbs. "carry" in "lakes carry" is "getupopa". "watch" in "hunters watch" is "sukupo".} \\
\texttt{{-} Let's verify roots for other nouns. "farmer" is "ti" in "big farmer" (Ex 1 object). In "five farmers" (Ex 3 object), it is "tigo lofa". Wait, Ex 3: "gugepopa ti totu tigo lofa". Object is "five farmers". So "farmer" base might be "ti"? But "ti" is also "road"? That can't be right.} \\
\texttt{{-} Let's re{-}read Ex 1: "the five lakes carry the big farmer". Subject: "the five lakes". Object: "the big farmer".} \\
\texttt{{-} If Object is "the big farmer", and output ends with "ti", then "ti" could be "farmer" or "the big farmer" or something else.} \\
\texttt{{-} Let's look at Ex 3: "a blue road feed the five farmers". Subject: "a blue road". Object: "the five farmers".} \\
\texttt{{-} Output: "gugepopa ti totu tigo lofa".} \\
\texttt{{-} If word order is SVO, Subject is "gugepopa ti", Verb is "totu", Object is "tigo lofa".} \\
\texttt{{-} So "road" with "blue" is "ti" (if "blue" is a prefix/suffix on the noun).} \\
\texttt{{-} In Ex 1, "the five lakes" is "gugepa fugo fogupo". So "lake" is likely "fugo". "five" is "fug{-}"? "the" is "{-}pa" or "{-}po"?} \\
\texttt{{-} Let's compare "lakes" (Ex 1) and "ships" (Ex 2: "kutopi").} \\
\texttt{{-} Ex 2: "five hunters watch the two red ships". Subject: "five hunters" {-}> "gisapopa fago te". Object: "the two red ships" {-}> "ti".} \\
\texttt{{-} Wait, in Ex 2, the object is "ti"? That matches the "farmer" object in Ex 1? That seems like a collision unless "farmer" and "ships" are the same word? Unlikely.} \\
\texttt{{-} Let's re{-}examine Ex 2 carefully. "five hunters watch the two red ships".} \\
\texttt{{-} Maybe the word order is not simple SVO? Or maybe I am misidentifying which part is subject/object.} \\
\texttt{{-} Input: `five hunters watch the two red ships`.} \\
\texttt{{-} Output: `gisapopa fago te kutopi sukupo ti`.} \\
\texttt{{-} If SVO: Subject=`gisapopa fago te`, Verb=`sukupo`, Object=`ti`.} \\
\texttt{{-} If VSO: Verb=`gisapopa fago te`, Subject=`sukupo`, Object=`ti`.} \\
\texttt{{-} If OVS: ...} \\
\texttt{{-} Let's look at Ex 1: "the five lakes carry the big farmer". Output: `gugepa fugo} \\
            }
        }
    }
    & & \\ \\

    \theutterance \stepcounter{utterance}  
    & & & \multicolumn{2}{p{0.3\linewidth}}{
        \cellcolor[rgb]{0.9,0.9,0.9}{
            \makecell[{{p{\linewidth}}}]{
                \texttt{\tiny{[GM$|$GM]}}
                \texttt{could not find any probe sentences; put them under a 'PROBES:' line, one per line} \\
            }
        }
    }
    & & \\ \\

\end{supertabular}
}

\end{document}
